\begin{frame}{Validation}{Overview}
  \begin{itemize}
    \item 1.1 Structural trade-off behaviour (threshold vs.\ time)
    \item 1.2 Cross-model scaling of emissions and savings
    \item 1.3 Regional and temporal realism
    \item 1.4 Checkpoint granularity and overhead effects
  \end{itemize}
  \vspace{0.5cm}
  \begin{block}{Purpose}
    Validation checks whether the simulation behaves plausibly for
    carbon-aware LLM
    training, i.e.\ whether it captures the expected trade-offs and patterns
    seen in real grids and training workloads, rather than just being
    ``correct code''.
  \end{block}
\end{frame}

% ---------------------------------------------------------------------------
% Detailed example: 1.3 Regional and Temporal Realism
% ---------------------------------------------------------------------------

% \begin{frame}{Validation}{1.3 Regional and Temporal Realism}
%   \textbf{Goal:} Check that regional and temporal variation in grid
%   carbon intensity
%   leads to plausible differences in baseline emissions and savings.
%   \vspace{0.3cm}

%   \textbf{Key inputs}
%   \begin{itemize}
%     \item Carbon intensity time series $\mathit{co2\_intensity}(t)$
%       per region and year.
%     \item Region and year choices (e.g.\ $\mathit{regions\_set}$,
%       $\mathit{historical\_years}$).
%     \item Fixed training profile (model, GPU count, power, training duration).
%   \end{itemize}
% \end{frame}

% \begin{frame}{Validation}{1.3 Regional and Temporal Realism}
%   \textbf{Expected behaviour}
%   \begin{itemize}
%     \item Define the average intensity for a region and year:
%       \[
%         \bar{c}_{\text{region,year}} =
%         \frac{1}{T} \int_0^T
%         \mathit{co2\_intensity}^{\text{region,year}}(t)\, dt.
%       \]
%     \item For fixed energy, baseline emissions should roughly follow
%       $\bar{c}_{\text{region,year}}$ (cleaner grids $\Rightarrow$
%       lower emissions).
%     \item Regions / periods with high temporal variability should offer
%       \emph{more} savings from pausing.
%   \end{itemize}
% \end{frame}

% \begin{frame}{Validation}{1.3 Regional and Temporal Realism}
%   \textbf{Validation procedure (1/2)}
%   \begin{enumerate}
%     \item For each region and year:
%       \begin{itemize}
%         \item Compute $\bar{c}_{\text{region,year}}$ from the raw time series.
%         \item Run a baseline simulation (no pausing) and obtain
%           $E_{\text{baseline}}$.
%       \end{itemize}
%     \item Check that the ordering of $E_{\text{baseline}}$ across regions/years
%       matches the ordering of $\bar{c}_{\text{region,year}}$
%       (allowing for small deviations).
%   \end{enumerate}
% \end{frame}

% \begin{frame}{Validation}{1.3 Regional and Temporal Realism}
%   \textbf{Validation procedure (2/2)}
%   \begin{enumerate}
%       \setcounter{enumi}{2}
%     \item Fix a carbon-aware pausing policy
%       (e.g.\ threshold, hysteresis, pause granularity).
%     \item Sweep over:
%       \begin{itemize}
%         \item Different regions.
%         \item Different start times within the year (seasonality).
%       \end{itemize}
%     \item Compare resulting
%       $\text{CO2\_savings\_pct}$ and $\text{time\_overhead\_pct}$:
%       \begin{itemize}
%         \item High-variance regions $\Rightarrow$ higher potential savings.
%         \item Low-variance, already-clean regions $\Rightarrow$
%           limited extra savings.
%       \end{itemize}
%   \end{enumerate}

%   \vspace{0.2cm}
%   This provides evidence that the simulation responds realistically
%   to regional and
%   temporal grid characteristics.
% \end{frame}
